\section{Two reflection-reduced closing equations}

Set
\begin{equation}\label{eq:chain}
 b=\frac{1-6a^2}{2},\qquad
 c=1-6b^2-a,\qquad
 d=1-6c^2-b.
\end{equation}
These are successive uses of \eqref{eq:recurrence} starting at a reflection
axis.

\subsection{The six-cycle \texorpdfstring{\(A_6\)}{A6}}

The vertex--vertex reflection pattern is
\begin{equation}\label{eq:A6pattern}
 (q_0,\ldots,q_5)=(a,b,c,d,c,b).
\end{equation}
The remaining endpoint equation is \(1-6d^2-2c=0\).  Its exact
factorization contains three quadratic factors, one quartic factor, and the
primitive factor
\begin{equation}\label{eq:A6coordinate}
 f_6(a)=2916a^6-1782a^4+108a^3+279a^2-33a-2.
\end{equation}
Reduction of all six recurrence residuals modulo \(f_6\) gives zero.  The
physical coordinate is the unique root in
\[
 \left(\frac{551907131}{10^9},\frac{551907132}{10^9}\right),
\]
on which the chain signs are \((+,-,-,-)\).  The full cyclic orbit therefore
has exactly one positive coordinate, as required by \(A_6=000021\).

\subsection{The shared seven-cycle chain}

The edge-reflection pattern is
\begin{equation}\label{eq:P7pattern}
 (q_0,\ldots,q_6)=(a,b,c,d,d,c,b).
\end{equation}
Now the last equation is \(1-6d^2-c-d=0\), and
\begin{equation}\label{eq:P7factor}
 1-6d^2-c-d=-\frac12(6a^2+2a-1)f_7(a),
\end{equation}
where \(f_7\) is the degree-fourteen polynomial in
Appendix~\ref{sec:ledger}.  Every recurrence residual vanishes modulo
\(f_7\).

Two rational boxes select the two physical embeddings:
\begin{align*}
 B_7:&\quad
 a\in\left(-\frac{600956965}{10^9},-\frac{600956964}{10^9}\right),
 &\operatorname{sgn}(a,b,c,d)&=(-,-,-,+),\\
 A_7:&\quad
 a\in\left(\frac{551935742}{10^9},\frac{551935743}{10^9}\right),
 &\operatorname{sgn}(a,b,c,d)&=(+,-,-,-).
\end{align*}
In the full palindromic patterns these are, up to cyclic rotation, the
two-positive word \(B_7\) and the one-positive word \(A_7\).
